\documentclass[12pt]{article}
\usepackage[utf8]{inputenc}
\usepackage[T1]{fontenc}
\usepackage{hyperref}
\title{Maraghei Language Resources for K-12 Curriculum Development}
\author{Compiled by Claude Sonnet 5 (claude-sonnet-5) for Andr\'as Kornai}
\date{}
\begin{document}
\maketitle
\section{Introduction}
Maraghei (ISO 639-3 vmh, Glottocode mara1373) is a Tatic (Northwestern Iranian) variety. Wikipedia places it in Gilan Province, in the upper Rudbar (Alamut) area, and not near the city of Maragheh in East Azerbaijan, whose name it resembles. Sources are mostly encyclopedic and weakly cited. No teaching materials were found.
\section{Primary Educational Resources}
None found.
\section{Linguistic and Reference Materials}
\begin{itemize}
\item Wikipedia, Maraghei dialect: \url{https://en.wikipedia.org/wiki/Maraghei_dialect}
\item Wikipedia, Maraghei: \url{https://en.wikipedia.org/wiki/Maraghei}
\item Verbix endangered-languages entry (status listed as dormant): \url{https://docs.verbix.com/EndangeredLanguages/Maraghei}
\end{itemize}
\section{Religious and Community Resources}
None found.
\section{Digital and Audio Resources}
None found.
\section{Curriculum Development Context}
There is essentially no published pedagogical material in the search results. Fieldwork and documentation would be needed before any curriculum could be built. The relationship to Talysh and Gozarkhani suggests that materials for related Tatic varieties may be worth searching for.
\end{document}